\documentclass[11pt,a4paper]{article}

% ==========================================
% PREÁMBULO Y PAQUETES
% ==========================================
\usepackage[utf8]{inputenc}
\usepackage[T1]{fontenc}
\usepackage[spanish, es-tabla]{babel}
\usepackage[margin=2.5cm]{geometry}
\usepackage{amsmath, amssymb}
\usepackage{booktabs}
\usepackage{graphicx}
\usepackage{hyperref}
\usepackage{enumitem}
\usepackage{multicol}
\usepackage{xcolor}
\usepackage{tikz}
\usetikzlibrary{arrows.meta, positioning, intersections, calc}
\usepackage[most]{tcolorbox}
\usepackage{titlesec}
\usepackage{float}
\usepackage{array}
\usepackage{longtable}
\usepackage{setspace}
\usepackage{lmodern}
\usepackage{fancyhdr}
\usepackage{lastpage}

% ==========================================
% CONFIGURACIÓN DE HIPERVÍNCULOS
% ==========================================
\hypersetup{
    colorlinks=true,
    linkcolor=blue!70!black,
    urlcolor=blue,
    pdftitle={Apunte EAE105A - Introducción a la Economía},
    pdfauthor={Juan Francisco Fuentes},
    pdfsubject={Apunte detallado de estudio},
    bookmarksnumbered=true
}

% ==========================================
% DEFINICIÓN DE COLORES
% ==========================================
\definecolor{azul}{RGB}{28, 84, 142}
\definecolor{azulosc}{RGB}{15, 50, 90}
\definecolor{verde}{RGB}{39, 128, 70}
\definecolor{rojo}{RGB}{158, 42, 42}
\definecolor{naranja}{RGB}{206, 118, 35}
\definecolor{gris}{RGB}{245, 245, 245}
\definecolor{morado}{RGB}{99, 69, 164}
\definecolor{amarillo}{RGB}{255, 193, 7}

% ==========================================
% ESTILOS DE SECCIONES
% ==========================================
\titleformat{\section}
  {\normalfont\Large\bfseries\color{azulosc}}
  {\thesection}{0.8em}{}[{\color{azul}\titlerule[0.8pt]}]
  
\titleformat{\subsection}
  {\normalfont\large\bfseries\color{azul}}
  {\thesubsection}{0.8em}{}

\titleformat{\subsubsection}
  {\normalfont\bfseries\color{azul!70!black}}
  {\thesubsubsection}{0.8em}{}

% ==========================================
% DEFINICIÓN DE CAJAS PEDAGÓGICAS
% ==========================================

% Caja: Concepto Clave
\newtcolorbox{concepto}[1]{
    colback=blue!5!white,
    colframe=azul,
    title=\textbf{🔑 Concepto clave:} #1,
    fonttitle=\bfseries\large,
    boxrule=1pt,
    arc=3mm,
    left=4mm, right=4mm, top=3mm, bottom=3mm,
    shadow={2mm}{-1mm}{0mm}{black!25},
    breakable
}

% Caja: Definición Formal
\newtcolorbox{definicion}[1]{
    colback=morado!5!white,
    colframe=morado,
    title=\textbf{📋 Definición:} #1,
    fonttitle=\bfseries,
    boxrule=0.8pt,
    arc=2mm,
    left=3mm, right=3mm, top=2mm, bottom=2mm,
    breakable
}

% Caja: Observación
\newtcolorbox{observacion}[1]{
    colback=yellow!8!white,
    colframe=naranja,
    title=\textbf{💡 Observación:} #1,
    fonttitle=\bfseries,
    boxrule=0.8pt,
    arc=2mm,
    left=3mm, right=3mm, top=2mm, bottom=2mm,
    breakable
}

% Caja: Advertencia
\newtcolorbox{advertencia}[1]{
    colback=red!5!white,
    colframe=rojo,
    title=\textbf{⚠️ ¡Ojo! (Trampa de prueba):} #1,
    fonttitle=\bfseries,
    boxrule=1pt,
    arc=2mm,
    left=3mm, right=3mm, top=2mm, bottom=2mm,
    borderline west={4pt}{0pt}{rojo},
    breakable
}

% Caja: Ejercicio Resuelto
\newtcolorbox{ejercicio}[1]{
    colback=green!5!white,
    colframe=verde,
    title=\textbf{✓ Ejercicio resuelto:} #1,
    fonttitle=\bfseries,
    boxrule=0.8pt,
    arc=2mm,
    left=3mm, right=3mm, top=2mm, bottom=2mm,
    breakable
}

% Caja: Ejemplo
\newtcolorbox{ejemplo}[1]{
    colback=yellow!3!white,
    colframe=amarillo!80!orange,
    title=\textbf{📌 Ejemplo:} #1,
    fonttitle=\bfseries,
    boxrule=0.8pt,
    arc=2mm,
    left=3mm, right=3mm, top=2mm, bottom=2mm,
    breakable
}

% Caja: Resumen
\newtcolorbox{resumen}[1]{
    colback=gray!8!white,
    colframe=black!40,
    title=\textbf{📌 Resumen:} #1,
    fonttitle=\bfseries,
    boxrule=0.8pt,
    arc=2mm,
    left=3mm, right=3mm, top=2mm, bottom=2mm,
    breakable
}

% Caja: Prueba Asociada
\newtcolorbox{prueba}[1]{
    colback=orange!8!white,
    colframe=naranja,
    title=\textbf{📝 #1},
    fonttitle=\bfseries\large,
    boxrule=1pt,
    arc=2mm,
    left=3mm, right=3mm, top=2mm, bottom=2mm,
    breakable
}

% ==========================================
% ESTILOS PERSONALIZADOS
% ==========================================
\newcommand{\nota}[1]{\textit{\color{azul}#1}}
\newcommand{\formula}[1]{\textbf{\color{rojo}#1}}
\newcommand{\clave}[1]{\textbf{\color{azulosc}#1}}
\newcommand{\respuesta}[1]{\textcolor{verde}{$\checkmark$ #1}}

% ==========================================
% ENCABEZADO Y PIE
% ==========================================
\pagestyle{fancy}
\fancyhf{}
\fancyhead[L]{\textbf{EAE105A - Introducción a la Economía}}
\fancyhead[R]{\thepage}
\fancyfoot[C]{\textit{Apunte detallado de estudio - Pontificia Universidad Católica}}
\renewcommand{\headrulewidth}{0.8pt}
\renewcommand{\footrulewidth}{0.8pt}

\setlength{\parindent}{0pt}
\setlength{\parskip}{0.5em}

% ==========================================
% DOCUMENTO COMIENZA AQUÍ
% ==========================================
\begin{document}

% ==========================================
% PORTADA
% ==========================================
\begin{titlepage}
    \centering
    \vspace*{2cm}
    
    {\scshape\Large Pontificia Universidad Católica de Chile \par}
    \vspace{0.5cm}
    {\scshape\large Instituto de Economía \par}
    
    \vspace{3cm}
    \rule{\textwidth}{2pt}
    \vspace{0.7cm}
    
    {\huge\bfseries Introducción a la Economía \par}
    {\LARGE\bfseries EAE105A \par}
    \vspace{0.5cm}
    {\Large\itshape El apunte definitivo: completo, detallado y listo para estudiar \par}
    
    \vspace{0.7cm}
    \rule{\textwidth}{2pt}
    \vspace{3cm}
    
    {\large \textbf{Autor:} Juan Francisco Fuentes Gonzáles \par}
    \vspace{1cm}
    {\large \textit{Material exhaustivo basado en:\par}}
    {\normalsize Clases, ayudantías, apuntes y pruebas anteriores \par}
    \vspace{2cm}
    
    {\large \textbf{Estructura del apunte:} \par}
    \begin{itemize}
        \item \textbf{Prueba 1:} Introducción + Conceptos generales + FPP + Comercio inicial
        \item \textbf{Prueba 2:} Microeconomía completa + Mercados + Demanda/Oferta + Equilibrio + Elasticidad + Bienestar
        \item \textbf{Prueba 3:} Comercio internacional + Monopolio + Externalidades + Bienes públicos + Macroeconomía
        \item \textbf{Examen:} Toda la materia con énfasis en temas clave no evaluados en pruebas
    \end{itemize}
    
    \vfill
    {\large \today \par}
    \vspace{0.5cm}
    {\small Compilado para Overleaf \par}
\end{titlepage}

\newpage

% ==========================================
% TABLA DE CONTENIDOS
% ==========================================
\tableofcontents

\newpage

% ==========================================
% MAPA VISUAL DEL CURSO
% ==========================================
\section*{Guía de lectura: ¿Qué estudiar para cada prueba?}
\addcontentsline{toc}{section}{Guía de lectura}

\begin{center}
\textbf{\Large ¿DÓNDE ENCONTRAR TUS TEMAS?}
\end{center}

\vspace{1cm}

\begin{prueba}{PRUEBA 1 (Introducción y Conceptos Base)}
\textbf{Temas:}
\begin{itemize}
    \item ✓ Motivación: por qué estudiar economía
    \item ✓ Conceptos generales (escasez, eficiencia, incentivos)
    \item ✓ Economía como ciencia (modelos, ceteris paribus, correlación vs causalidad)
    \item ✓ Aplicación al comercio: FPP, ventaja absoluta y comparativa, autarquía y comercio
\end{itemize}
\textbf{Secciones del apunte:} Parte I, Secciones 1-4
\end{prueba}

\vspace{0.7cm}

\begin{prueba}{PRUEBA 2 (Microeconomía y Mercados)}
\textbf{Temas:}
\begin{itemize}
    \item ✓ Mercados y competencia
    \item ✓ Teoría de la demanda y oferta
    \item ✓ Equilibrio de mercado y estática comparativa
    \item ✓ Elasticidad (precio de la demanda y oferta)
    \item ✓ Eficiencia del mercado (excedentes)
    \item ✓ Fijación de precios (máximos y mínimos)
    \item ✓ Impuestos y subsidios
\end{itemize}
\textbf{Secciones del apunte:} Parte II, Secciones 1-7
\end{prueba}

\vspace{0.7cm}

\begin{prueba}{PRUEBA 3 (Comercio, Monopolio, Fallas y Macroeconomía)}
\textbf{Temas:}
\begin{itemize}
    \item ✓ Comercio internacional (exportaciones, importaciones, aranceles)
    \item ✓ Monopolio
    \item ✓ Externalidades
    \item ✓ Bienes públicos y recursos comunes
    \item ✓ Contabilidad nacional y PIB
    \item ✓ Inflación y IPC
    \item ✓ Sistema monetario
\end{itemize}
\textbf{Secciones del apunte:} Parte II (Sección 8) + Parte III (Secciones 1-4)
\end{prueba}

\vspace{0.7cm}

\begin{prueba}{EXAMEN (Toda la materia)}
\textbf{Con énfasis en:}
\begin{itemize}
    \item Conceptos clave de todas las pruebas
    \item Conexiones entre temas
    \item Desempleo, política fiscal y deuda pública
    \item Aplicaciones a situaciones reales
\end{itemize}
\textbf{Secciones del apunte:} TODO (Partes I, II y III)
\end{prueba}

\newpage

\end{document}
